\documentclass[10pt,a4paper]{amsart}
\usepackage[margin=19mm]{geometry}
\usepackage{amsmath,amssymb,mathtools}
\usepackage[scr=rsfs,cal=euler]{mathalfa}
\usepackage{xfrac}
\usepackage[unicode,hidelinks,bookmarks=false]{hyperref}
\newcommand{\ie}{i.e.\ }
\newcommand{\cf}{cf.\ }
\newcommand{\resp}{respectively\ }
\newcommand{\Subsection}{\subsection}
\providecommand{\nobreakdash}{\nobreak\hbox{-}}
\setlength{\emergencystretch}{2em}
\multlinegap0pt
\usepackage{xcolor}
\usepackage{mathtools}
\usepackage{amssymb}
\usepackage{enumerate}
\usepackage[shortlabels]{enumitem}
\usepackage{float}
\usepackage{tikz}
\newtheorem{thm}{Theorem}
\title{A counter-example to persistence in generalised preferential attachment trees}
\author{Tejas Iyer}
\date{Source: arXiv:2604.15007v1 (CC BY 4.0)}
\begin{document}
\maketitle

\noindent\emph{Source and licence.} A counter-example to persistence in generalised preferential attachment trees. Version arXiv:2604.15007v1 (CC BY 4.0), distributed by arXiv under the Creative Commons Attribution 4.0 International licence, http://creativecommons.org/licenses/by/4.0/, as recorded in the arXiv licence metadata for this exact version and shown on the abstract page https://arxiv.org/abs/2604.15007v1. Full licence notice: \texttt{licenses/arxiv-2604.15007-CC-BY-4.0.txt}.\par\smallskip

\noindent\textit{Editorial scope.} Counted unit: the article's thm thm:counter-example (source0.tex lines 135--137). The article's complete original proof of it is delivered with it (source0.tex lines 154--168, one \texttt{proof} environment of 15 lines, which the article places away from the statement). No mathematics is written, repaired or re-derived: the statement and the proof are the article's own lines, in the article's own notation, and no retained line is substituted or re-flowed.  No further source-local context is needed: every cross-reference of every retained line resolves inside the two delivered spans.  Nothing was omitted from any retained span, so the package carries no omission record.  The retained lines cite 1 works, whose entries are transcribed verbatim from the article's own bibliography source in the e-print and delivered with short editorial labels recorded in the item's reference mapping.  No retained line contains a figure environment, an image-inclusion command or drawing code, so no image is delivered and the package declares no asset dependency.  Editorial formatting only, in addition: an AMS \texttt{amsart} preamble that restates the article's own declarations and macros used by the retained lines, namely the environments \texttt{thm} on separate counters, together with the standard packages those lines need.  The retained environments print in this document as Theorem 1 in order of appearance, whereas the article prints the same items with the numbering of its own full document.  The complete native source of the article as distributed in the arXiv e-print for this exact version is bundled unchanged as \texttt{sources/198574/source0.tex}.
\begin{thm} \label{thm:counter-example}
    There exists a function $f\colon \mathbb{N}_0 \rightarrow (0,\infty)$ satisfying $\sum_{j=0}^{\infty} \frac{1}{f(j)^2} < \infty$, such that, almost surely, for the associated generalised preferential attachment trees $(\mathcal{T}_{n})_{n \in \mathbb{N}_0}$ there is no persistent hub. 
\end{thm}

\begin{proof}[Proof of Theorem~\ref{thm:counter-example}]
    Note that if $(f(j))_{j \in \mathbb{N}_0}$ is a sequence such that $\sum_{j=0}^{\infty} \frac{1}{f(j)} = \infty$, and yet \[\lim_{n \to \infty} \tau_{n} =: \tau_{\infty} < \infty\] almost surely, then it is impossible for $(\mathcal{T}_{n})_{n \in \mathbb{N}_0}$ to contain a persistent hub. Indeed, $\sum_{j=0}^{\infty} \frac{1}{f(j)} = \infty$ ensures, by standard criteria for pure-birth processes, that $\xi([0,t]) < \infty$ for all $t \in [0, \infty)$ which means, by a union bound over the countable set $\mathcal{U}$, that the degree of every node in $\mathcal{T}_{\infty}:= \bigcup_{n \in \mathbb{N}_0} \mathcal{T}_n$ is finite almost surely. If the maximal degree in $\mathcal{T}_{\infty}$ was bounded from above, by $L$ say, this would imply the existence of infinitely many individuals whose ancestors all have at most $L$ children, which contradicts~\cite[Proposition~4.4]{inhom-sup-pref-attach} (such individuals are termed $L$-moderate there). Therefore, the maximal degree of $\mathcal{T}_{n}$, and hence the degree of any persistent hub must diverge to infinity.  
    
    We now make a minor modification to the construction of Galganov and Ilienko, to construct a function $f$ such that
    $\sum_{j=0}^{\infty} \frac{1}{f(j)} = \infty$, $\sum_{j=0}^{\infty} \frac{1}{f(j)^2} < \infty$ and $\tau_{\infty} < \infty$ almost surely. To do so, we define the values $f(0), f(1), \ldots$, by 
    \[
    1, \underbrace{\alpha_{1}, \cdots, \alpha_1}_{d_1 \text{ times}}, 2, \underbrace{\alpha_{2}, \cdots, \alpha_2}_{d_2 \text{ times}}, 3, \underbrace{\alpha_{2}, \cdots, \alpha_3, }_{d_3 \text{ times}} \cdots, 
    \]
    where $d_{k} := 4^{k^2}$ and $\alpha_{k} := 2^{2^{k^3}}$. Note that $\sum_{j=0}^{\infty} \frac{1}{f(j)} > \sum_{j=1}^{\infty} \frac{1}{j} = \infty$, whilst $\sum_{j=0}^{\infty} \frac{1}{f(j)^{2}} = \sum_{k=1}^{\infty} \frac{1}{k^2} + \sum_{k=1}^{\infty} \frac{4^{k^2}}{2^{2^{k^3}}} < \infty$. 
    It remains to show that $\tau_{\infty} < \infty$ almost surely. For this, note that Galganov and Ilienko have already shown that this occurs for the process with smaller rates 
    \[
    1, \underbrace{\alpha_{1}, \cdots, \alpha_1}_{d_1 \text{ times}}, 1, \underbrace{\alpha_{2}, \cdots, \alpha_2}_{d_2 \text{ times}}, 1, \underbrace{\alpha_{3}, \cdots, \alpha_3, }_{d_3 \text{ times}} \cdots, 
    \]
    whence, by a standard coupling argument (coupling the pure-birth processes associated with individuals in the processes), we also get $\tau_{\infty} < \infty$ in this process. 
\end{proof}

\begin{thebibliography}{99}
\bibitem[inhom-]{inhom-sup-pref-attach} T.~Iyer and B.~Lodewijks. \newblock On the structure of genealogical trees associated with explosive {Crump-Mode-Jagers} branching processes. \newblock arXiv preprint arXiv:2311:14664, 2023.
\end{thebibliography}
\end{document}
